\section{Sikhism}

Sikhism is the tradition founded by Guru Nanak and the nine Gurus who followed
him, with its scripture the Guru Granth Sahib. It teaches a strict oneness. One
formless creator is present in all things and beyond all things, and the soul's
task is to remember that one and merge back into it. It rejects idol worship and
the divisions of caste, holding that the one light dwells equally in all.

\subsection{The Creation Model}

For countless ages there was only formless dark. No earth, no sky, no day or
night, no measure of time. There was only the one, alone and unmanifest, held
in its own will. Then the one willed it, and with a single word brought forth
the worlds. From that one utterance came the air, and from the air the water,
and from the water the play of living things. The one set the worlds turning by
its own order and poured itself into all of them. It is not a maker standing
outside its work. It is present in every creature and every grain, the formless
one within all forms. The same will that brought the worlds out holds them in
being and will draw them home again. The whole of life is the unfolding of that
one will, and the goal of the soul is to lose its separateness and merge back
into the one from which it came.

\subsection{The One}

The opening words of the Sikh scripture are \textbf{Ik Onkar}, \enquote{there is
one.} It is the root statement of the whole faith. There is one supreme being,
one creator, immanent in all and beyond all. From it follows the
opening verse (the \textbf{Mul Mantar}) that describes the one through a short
chain of attributes.

\begin{longtable}{L{0.24\linewidth} L{0.58\linewidth}}
\caption{The Mul Mantar, the description of the one.}\\
\toprule
Term & Meaning \\
\midrule
Ik Onkar & There is one supreme being \\
Satnam & Truth is its name \\
Karta Purakh & The creator being \\
Nirbhau & Without fear \\
Nirvair & Without enmity \\
Akal Murat & Timeless in form \\
Ajuni & Beyond birth and death \\
Saibhang & Self-existent \\
Gurprasad & Known by the grace of the teacher \\
\bottomrule
\end{longtable}

The one is timeless and self-existent, without fear and without hatred, never
born and never dying. The devotional name for this one is the \enquote{wondrous
lord} (\textbf{Waheguru}). Sikhism holds a strict oneness. It rejects the
worship of idols and images, and it rejects the divisions of caste, since the
one light dwells equally in all.

\subsection{The Name and the Order}

Two words carry the inner life of the faith. \textbf{Naam} is the divine name,
the presence of the one made available to the heart. \textbf{Naam simran} is
the practice of remembrance, holding the name in mind until the heart itself
remembers. Remembering the name cleanses the self and draws it toward the one.

\textbf{Hukam} is the divine order, the will by which all things happen. By the
hukam life rises and falls, joy and sorrow come, and the soul is bound to the
cycle of birth or released from it. To live well is to read the hukam and act
in step with it rather than against it.

\subsection{The Thieves and the Virtues}

What blocks the soul is egoism or \enquote{I-am-ness} (\textbf{haumai}), the
sense of a separate self standing apart from the one. Haumai is the root of the
\textbf{five thieves}, the impulses that rob the soul and keep it far from the
one. Against them stand the \textbf{five virtues}.

\begin{longtable}{L{0.36\linewidth} L{0.36\linewidth}}
\caption{The five thieves and the five virtues.}\\
\toprule
Five thieves & Five virtues \\
\midrule
Kaam (lust) & Sat (truth) \\
Krodh (wrath) & Daya (compassion) \\
Lobh (greed) & Santokh (contentment) \\
Moh (attachment) & Nimrata (humility) \\
Ahankar (pride) & Pyaar (love) \\
\bottomrule
\end{longtable}

\subsection{The Five Realms}

The Japji describes the ascent of the soul as a passage through five
realms or stages (\textbf{khands}). They rise from outward duty to final union
with the one.

\begin{longtable}{L{0.06\linewidth} L{0.22\linewidth} L{0.52\linewidth}}
\caption{The five khands of the Japji, the stages of the soul's ascent.}\\
\toprule
\# & Realm & Meaning \\
\midrule
1 & Dharam Khand & The realm of duty, righteous action and law \\
2 & Gian Khand & The realm of knowledge, the vastness of creation \\
3 & Saram Khand & The realm of effort, the shaping of the spirit \\
4 & Karam Khand & The realm of grace, where the one acts in the soul \\
5 & Sach Khand & The realm of truth, union with the one \\
\bottomrule
\end{longtable}

The first realm sets the soul in right action. The second opens its eyes to the
breadth of creation. The third is the realm of spiritual effort, where the inner
self is reshaped. The fourth is grace, no longer effort but gift. The fifth is
the realm of truth (\textbf{Sach Khand}), where the soul attains union with the
one. The end of the path is liberation (\textbf{mukti}), the loss of haumai and
the merging of the soul back into the one.

\subsection{Comparison to the Base Models}

Sikhism gives the most compressed form of the shared pattern. One formless
source wills the worlds into being with a single word, pours itself into all of
them, and draws the soul back through ranked stages to merge into the one. The
five khands are a clear version of the staged ascent seen across the base
models, and mukti is the go-out-and-return stated as the dissolving of the
separate self. What sets it apart is its strict oneness and its rejection of
idol and caste. The same one that goes out as the worlds is the one the soul
returns to.
